\section{The relative strengthening in context}\label{sec:comparison}

The main theorem is not a new proof that planar convex codes are polygonal;
that result is due to \citet{BukhJeffs2023}.  Its contribution is to make a
finite collection of geometric data survive the same shrinking process.
The distinction is summarized in \cref{tab:comparison}.

\begin{table}[h]
\centering
\small
\begin{tabular}{@{}p{.25\textwidth}p{.29\textwidth}p{.35\textwidth}@{}}
\toprule
Feature & Planar polygonalization & Present relative theorem\\
\midrule
Carrier & An open planar realization & A fixed compact polygon and its relative interior\\
Exact neural code & Preserved & Preserved relative to the fixed carrier\\
Prescribed lines & Used internally for closure-only patterns & Arbitrary finite lines added to the internal family\\
Strict traces & No external trace constraint in the theorem statement & Exact source subsets retained after clipping\\
Closed contacts & Used as proof data where needed & Any prescribed point or segment section may be retained\\
Atom witnesses & A representative of each codeword is used & Arbitrarily designated witnesses remain at their original coordinates\\
Interface consequence & Not part of the theorem statement & Shared edges and finite facewise models agree pointwise\\
Higher-dimensional filling & Not addressed & Explicitly not claimed\\
\bottomrule
\end{tabular}
\caption{Comparison with the baseline planar theorem.  The present result
strengthens its finite constraints, not its planar dimension.}
\label{tab:comparison}
\end{table}

The proof does not approximate the fields in Hausdorff distance.  It shrinks
the closed tuple while preserving a finite representative package.  The
trace frames, ambient vertices, and protected triangles become additional
immutable data in the inclusion-minimal construction.  This is why exact
atoms and exact line sections can be controlled simultaneously even in the
presence of tangencies and closure-supported behavior.
